\section*{Chapter \ref{ch:g9:maps-of-the-earth} --- Maps of the Whole Earth}

\begin{solution}{exo:g9:maps-of-the-earth:1}
(a) A \omterm{def:g9:maps-of-the-earth:projection}{map projection}. (b) A \omterm{def:g9:maps-of-the-earth:contour}{contour line}. (c) A \omterm{def:g9:maps-of-the-earth:numerical-scale}{numerical scale}. (d) The
\omterm{def:g9:maps-of-the-earth:contour}{contour interval}. (e) \omterm{def:g9:maps-of-the-earth:remote-sensing}{Remote sensing}.
\end{solution}

\begin{solution}{exo:g9:maps-of-the-earth:2}
(a) $4 \times 25\,000 = \num{100000}$~cm $= \qty{1}{km}$. (b) $3 \times
50\,000 = \num{150000}$~cm $= \qty{1.5}{km}$. (c) $2.5 \times 1\,000\,000
= \num{2500000}$~cm $= \qty{25}{km}$. (d) $\qty{1}{km} = \num{100000}$~cm;
$100\,000 \div 25\,000 = \qty{4}{cm}$.
\end{solution}

\begin{solution}{exo:g9:maps-of-the-earth:3}
(a) False: a sphere cannot be flattened without distorting lengths, so
every \omterm{def:g3:continents-and-oceans:world-map}{world map} distorts something. (b) True. (c) False: an equal-area \omterm{def:g2:drawing-to-map:map}{map}
keeps areas, but bends shapes, especially near its edges. (d) True.
\end{solution}

\begin{solution}{exo:g9:maps-of-the-earth:4}
(a) The slope is steep. (b) The ground is gentle. (c) A \omterm{def:g1:landscapes-around-us:landform}{hill}, its top
inside the smallest line. (d) A point of the ground has only one \omterm{def:g4:world-physical-map:altitude}{altitude},
so it cannot lie on two lines of different \omterm{def:g4:world-physical-map:altitude}{altitudes}.
\end{solution}

\begin{solution}{exo:g9:maps-of-the-earth:5}
(a) At \ang{60}\,N the \omterm{def:g1:landscapes-around-us:water}{lake} is drawn 4 times too large: it looks as large as
$4 \times 1000 = \qty{4000}{km^2}$. At the \omterm{def:g3:continents-and-oceans:equator}{equator} areas are not enlarged, so the
\omterm{def:g1:landscapes-around-us:water}{lake} drawn with the same size has a true area of \qty{4000}{km^2}.
(b) $4 \div 1.33 \approx 3$ times larger. (c) Areas, and distances north
to south, grow without limit towards the \omterm{def:g3:continents-and-oceans:poles}{poles}: the \omterm{def:g3:continents-and-oceans:poles}{poles} would be
infinitely far from the \omterm{def:g3:continents-and-oceans:equator}{equator}.
\end{solution}

\begin{solution}{exo:g9:maps-of-the-earth:6}
(a) \qty{50}{m}. (b) The summit is at \qty{544}{m}; the index contour around
it is the \qty{500}{m} line. (c) In the west: a \omterm{def:g1:landscapes-around-us:water}{river} flows downhill, so
its floor is lower downstream, to the east. (d) Between about the 300 and
the \qty{400}{m} lines: there the slope is steepest, gentler both near the
\omterm{def:g1:landscapes-around-us:landform}{valley} floor and near the summit.
\end{solution}

\begin{solution}{exo:g9:maps-of-the-earth:7}
(a) C at about \qty{400}{m}, D at about \qty{530}{m}. (b) About
\qty{135}{m}, about \qty{1.3}{km} from C, on the \omterm{def:g1:landscapes-around-us:landform}{valley} floor. (c) About
0.3 to \qty{0.4}{km}. (d) The heights are small compared with the
distances: without stretching, the \omterm{def:g1:landscapes-around-us:landform}{hills} would look almost flat.
\end{solution}

\begin{solution}{exo:g9:maps-of-the-earth:8}
(a) $2 \times 25\,000 = \num{50000}$~cm $= \qty{500}{m}$. (b) $300 - 200 =
\qty{100}{m}$. (c) $100 \div 500 = 0.2$, a slope of \qty{20}{\%}. (d) Gentler:
\qty{4}{cm} is \qty{1000}{m}, so the slope is $100 \div 1000 = \qty{10}{\%}$,
half as steep.
\end{solution}

\begin{solution}{exo:g9:maps-of-the-earth:9}
(a) $14 \times 18 = 252$ orbits. (b) $14 \times 103 = 1442$ minutes, about
\qty{24}{h}: a day. (c) $5.5 \div 1 = 5.5$ times. (d) For example: a \omterm{def:g8:ice-on-earth:glacier}{glacier}
that retreats, a forest cut down, a town that grows, a \omterm{def:g1:landscapes-around-us:water}{coast} or a \omterm{def:g1:landscapes-around-us:water}{river} that
moves.
\end{solution}

\begin{solution}{exo:g9:maps-of-the-earth:10}
(a) $20\,200 \div 300\,000 \approx 0.067$~s $= 6.7 \times 10^{-2}$~s.
(b) $300\,000 \times 10^{-3} = \qty{300}{km}$. (c) A tiny error of time
makes a huge error of distance: the clocks have to be extremely precise.
\end{solution}

\begin{solution}{exo:g9:maps-of-the-earth:11}
(a) Distance to A: $\sqrt{3.6^2 + 4.8^2} = \sqrt{12.96 + 23.04} =
\sqrt{36} = 6$; distance to B: $\sqrt{6.4^2 + 4.8^2} = \sqrt{40.96 +
23.04} = \sqrt{64} = 8$. (b) $(3.6, -4.8)$, the mirror image across the
line AB. (c) Distance from C to $(3.6, 4.8)$: $\sqrt{3.6^2 + 5.2^2} =
\sqrt{40} \approx 6.3$; to $(3.6, -4.8)$: $\sqrt{3.6^2 + 14.8^2} =
\sqrt{232} \approx 15.2$. The receiver is at $(3.6, 4.8)$. (d) One distance
gives a circle (a sphere in space), two give two possible points; another
distance is needed to choose.
\end{solution}

\begin{solution}{exo:g9:maps-of-the-earth:12}
(a) $29 \div 2.1 \approx 14$ times. (b) $2.1 \times 10.5 \approx
22$~million~\unit{km^2}. (c) Greenland seems to cover about 22~million
\unit{km^2}, close to Africa's 29~million: on the Mercator wall \omterm{def:g2:drawing-to-map:map}{map}, the two
look about the same size, although Africa is 14 times larger. (d) An
\omterm{def:g9:maps-of-the-earth:projection}{equal-area projection}.
\end{solution}

\begin{solution}{pb:g9:maps-of-the-earth:1}
\textbf{1.} $50\,000$~cm $= \qty{500}{m}$.

\textbf{2.} $3.6 \times 500 = \qty{1800}{m} = \qty{1.8}{km}$.

\textbf{3.} B--W: $2.5 \times 500 = \qty{1250}{m}$; W--T: $2.35 \times 500
\approx \qty{1180}{m}$.

\textbf{4.} Towards the north-east.

\textbf{5.} \qty{50}{m}.

\textbf{6.} S--B goes down: $296 - 143 = \qty{153}{m}$.

\textbf{7.} B--W: $335 - 143 = \qty{192}{m}$; W--T: $544 - 335 =
\qty{209}{m}$.

\textbf{8.} $209 \div 1180 \approx 0.18$, about \qty{18}{\%}.

\textbf{9.} A straight line from S to T would cross the \omterm{def:g1:landscapes-around-us:water}{river} where there
is no bridge; the path goes down to the bridge at B.

\textbf{10.} $1.8 + 1.25 + 1.18 \approx \qty{4.2}{km}$ at \qty{4}{km} an
hour: about 63 minutes; climb $192 + 209 = \qty{401}{m}$: about 40 minutes
more. In all, about 1~h~43~min.

\textbf{11.} From T she goes down to B, then climbs from B up to S: $296 -
143 = \qty{153}{m}$.

\textbf{12.} Length: $1.8 + 1.25 + 1.18 \approx \textbf{4.2 km}$; total
climb: $192 + 209 = \textbf{401 m}$.
\end{solution}
